\documentclass[12pt,a4paper]{article}

% Encoding and fonts
\usepackage[utf8]{inputenc}
\usepackage[T1]{fontenc}
\usepackage{lmodern}

% Page layout
\usepackage[top=1in, bottom=1in, left=1in, right=1in]{geometry}

% Typography
\usepackage{microtype}
\usepackage{parskip}

% Language
\usepackage[english]{babel}

% Headers and footers
\usepackage{fancyhdr}
\pagestyle{fancy}
\fancyhf{}
\fancyhead[L]{\leftmark}
\fancyhead[R]{\thepage}
\fancyfoot[C]{\thepage}
\renewcommand{\headrulewidth}{0.4pt}
\renewcommand{\footrulewidth}{0pt}

% Section styling
\usepackage{titlesec}
\titleformat{\section}{\Large\bfseries}{\thesection}{1em}{}
\titleformat{\subsection}{\large\bfseries}{\thesubsection}{1em}{}
\titleformat{\subsubsection}{\normalsize\bfseries}{\thesubsubsection}{1em}{}
\titlespacing*{\section}{0pt}{2.5ex plus 1ex minus .2ex}{1.5ex plus .2ex}
\titlespacing*{\subsection}{0pt}{2ex plus 1ex minus .2ex}{1ex plus .2ex}
\titlespacing*{\subsubsection}{0pt}{1.5ex plus 1ex minus .2ex}{0.8ex plus .2ex}

% List formatting
\usepackage{enumitem}
\setlist[itemize]{topsep=0.5ex, itemsep=0.25ex, parsep=0pt}
\setlist[enumerate]{topsep=0.5ex, itemsep=0.25ex, parsep=0pt}

% Essential packages
\usepackage{graphicx}
\usepackage[table]{xcolor}
\usepackage{booktabs}
\usepackage{array}
\usepackage{tabularx}
\usepackage{longtable}
\usepackage{amsmath}
\usepackage{amssymb}
\usepackage[normalem]{ulem}
\rowcolors{2}{gray!10}{white}
\renewcommand{\arraystretch}{1.3}

% Cross-references
\usepackage{hyperref}
\usepackage{cleveref}

% Hyperref setup
\hypersetup{
    colorlinks=true,
    linkcolor=blue,
    filecolor=magenta,
    urlcolor=cyan,
}

% Code listings
\usepackage{listings}
\definecolor{codebg}{RGB}{248,248,248}
\definecolor{codeframe}{RGB}{200,200,200}
\definecolor{codekeyword}{RGB}{0,0,180}
\definecolor{codecomment}{RGB}{0,128,0}
\definecolor{codestring}{RGB}{163,21,21}
\lstset{
    basicstyle=\ttfamily\small,
    breaklines=true,
    breakatwhitespace=false,
    frame=single,
    framerule=0.5pt,
    rulecolor=\color{codeframe},
    backgroundcolor=\color{codebg},
    keepspaces=true,
    showstringspaces=false,
    tabsize=4,
    xleftmargin=1em,
    xrightmargin=1em,
    aboveskip=1em,
    belowskip=1em,
    keywordstyle=\color{codekeyword}\bfseries,
    commentstyle=\color{codecomment}\itshape,
    stringstyle=\color{codestring},
}

% YAML language definition for listings
\lstdefinelanguage{YAML}{
    keywords={true,false,null,y,n},
    keywordstyle=\color{blue}\bfseries,
    sensitive=false,
    comment=[l]{\#},
    morecomment=[s]{/*}{*/},
    commentstyle=\color{gray}\ttfamily,
    stringstyle=\color{red}\ttfamily,
    morestring=[b]',
    morestring=[b]",
}

% TOML language definition for listings
\lstdefinelanguage{TOML}{
    keywords={true,false},
    keywordstyle=\color{blue}\bfseries,
    sensitive=true,
    comment=[l]{\#},
    commentstyle=\color{gray}\ttfamily,
    stringstyle=\color{red}\ttfamily,
    morestring=[b]',
    morestring=[b]",
}

% Dockerfile language definition for listings
\lstdefinelanguage{Dockerfile}{
    keywords={FROM,RUN,CMD,LABEL,EXPOSE,ENV,ADD,COPY,ENTRYPOINT,VOLUME,USER,WORKDIR,ARG,ONBUILD,STOPSIGNAL,HEALTHCHECK,SHELL},
    keywordstyle=\color{blue}\bfseries,
    sensitive=false,
    comment=[l]{\#},
    commentstyle=\color{gray}\ttfamily,
    stringstyle=\color{red}\ttfamily,
    morestring=[b]',
    morestring=[b]",
}

% JSON language definition for listings
\lstdefinelanguage{JSON}{
    keywords={true,false,null},
    keywordstyle=\color{blue}\bfseries,
    sensitive=true,
    commentstyle=\color{gray}\ttfamily,
    stringstyle=\color{red}\ttfamily,
    morestring=[b]",
}

% Code listing float environment
\usepackage{float}
\newfloat{listing}{htbp}{lol}[section]
\floatname{listing}{Listing}

% Admonition boxes
\usepackage{tcolorbox}
\tcbuselibrary{skins,breakable}

% Admonition style definitions
\newtcolorbox{admonitionbox}[2][]{
    enhanced,
    breakable,
    colback=#2!5,
    colframe=#2!80!black,
    fonttitle=\bfseries,
    title=#1,
    left=2mm,
    right=2mm,
}

% Synthon tuple boxes
\newtcolorbox{synthonbox}{
    enhanced,
    colback=blue!5,
    colframe=blue!75!black,
    fontupper=\large,
    center upper,
    boxrule=1pt,
    arc=4pt,
}

\begin{document}

\title{SynthOmnicon --- Lean 4 Formalization}
\date{\today}

\maketitle

Lean 4 / Mathlib formalizations connected to the SynthOmnicon framework.
Toolchain: \textbf{Lean 4.28.0} $\cdot$ \textbf{Mathlib v4.28.0}

\begin{center}
\rule{0.5\textwidth}{0.4pt}
\end{center}

\subsection{Library structure}

\begin{lstlisting}
SynthOmnicon/
  Primitives/
    Core.lean         -- 12 primitive types with lattice structure + cross-primitive axioms
    OPN_2adic.lean    -- Classical number theory: 2-adic valuation + Touchard congruence for OPNs
\end{lstlisting}

\begin{center}
\rule{0.5\textwidth}{0.4pt}
\end{center}

\subsection{\texttt{OPN\_2adic.lean} --- What is actually proven}

This file formalizes the \textbf{2-adic and 3-adic valuation arguments} constraining the structure of
odd perfect numbers (OPNs), culminating in a machine-verified proof of \textbf{Touchard's congruence
(1953)}. Every \texttt{sorry} is an honest marker of a genuine open problem; no result is overclaimed.

\subsubsection{Definitions}

\begin{table}[htbp]
\centering
\small
\rowcolors{1}{white}{white}
\begin{tabularx}{\textwidth}{>{\centering\arraybackslash}X >{\centering\arraybackslash}X}
\toprule
\rowcolor{gray!25}\textbf{Symbol} & \textbf{Definition} \\
\midrule
\rowcolors{1}{white}{gray!10}
\texttt{Perfect n} & \texttt{sigma(n) = 2n}, using Mathlib's \texttt{ArithmeticFunction.sigma 1} \\
\texttt{v₂ n} & \texttt{(Nat.factorization n) 2} --- 2-adic valuation via prime factorization \\
\bottomrule
\end{tabularx}
\end{table}

\begin{center}
\rule{0.5\textwidth}{0.4pt}
\end{center}

\subsubsection{Helper lemmas --- fully proved, no sorry}

\paragraph{2-adic modular arithmetic}

\begin{table}[htbp]
\centering
\small
\rowcolors{1}{white}{white}
\begin{tabularx}{\textwidth}{>{\centering\arraybackslash}X >{\centering\arraybackslash}X}
\toprule
\rowcolor{gray!25}\textbf{Lemma} & \textbf{Statement} \\
\midrule
\rowcolors{1}{white}{gray!10}
\texttt{pred\_dvd\_pow\_sub\_one} & \texttt{(p-1) ∣ (pⁿ-1)} --- via \texttt{geom\_sum\_mul} over ℤ \\
\texttt{v2\_div\_of\_dvd} & \texttt{v₂(a/b) = v₂(a) - v₂(b)} when \texttt{b ∣ a} \\
\texttt{v2\_eq\_one\_of\_mod4\_eq2} & \texttt{n \% 4 = 2 -\textgreater{} v₂(n) = 1} --- via \texttt{Prime.pow\_dvd\_iff\_le\_factorization} \\
\texttt{pow\_odd\_of\_odd} & \texttt{q \% 2 = 1 -\textgreater{} q\^{}i \% 2 = 1} \\
\texttt{pow\_mod4\_of\_mod4} & \texttt{p \% 4 = 1 -\textgreater{} p\^{}i \% 4 = 1} \\
\texttt{geom\_sum\_mod2} & \texttt{(\sum\_\{i\textless{}n\} q\^{}i) \% 2 = n \% 2} when \texttt{q} is odd \\
\texttt{geom\_sum\_odd\_mod2} & \texttt{(\sum\_\{i\textless{}2e+1\} q\^{}i) \% 2 = 1} --- odd count of odd terms is odd \\
\texttt{geom\_sum\_mod4} & \texttt{(\sum\_\{i\textless{}n\} p\^{}i) \% 4 = n \% 4} when \texttt{p \equiv 1 (mod 4)} \\
\texttt{geom\_sum\_mod4\_eq2} & Corollary: \texttt{n \% 4 = 2 -\textgreater{} (\sum\_\{i\textless{}n\} p\^{}i) \% 4 = 2} \\
\bottomrule
\end{tabularx}
\end{table}

\paragraph{$\sigma$ identities and product formula}

\begin{table}[htbp]
\centering
\small
\rowcolors{1}{white}{white}
\begin{tabularx}{\textwidth}{>{\centering\arraybackslash}X >{\centering\arraybackslash}X}
\toprule
\rowcolor{gray!25}\textbf{Lemma} & \textbf{Statement} \\
\midrule
\rowcolors{1}{white}{gray!10}
\texttt{sigma\_mul\_of\_coprime} & \texttt{sigma(ab) = sigma(a)sigma(b)} for \texttt{gcd(a,b)=1} --- from \texttt{isMultiplicative\_sigma} \\
\texttt{sigma\_prime\_pow\_ratio} & \texttt{sigma(pᵏ).(p-1) + 1 = p\^{}(k+1)} --- geometric sum identity in ℕ \\
\texttt{sigma\_prime\_pow\_lt} & \texttt{sigma(pᵏ).(p-1) \textless{} p\^{}(k+1)} --- the ratio \texttt{sigma(pᵏ)/pᵏ} never reaches the ceiling \texttt{p/(p-1)} \\
\texttt{opn\_mod4} & Any OPN \texttt{\equiv 1 (mod 4)} --- from \texttt{p \equiv 1 (mod 4)} and \texttt{m} odd $\rightarrow$ \texttt{m\^{}2 \equiv 1 (mod 4)} \\
\bottomrule
\end{tabularx}
\end{table}

\paragraph{3-adic modular arithmetic (for Touchard)}

\begin{table}[htbp]
\centering
\small
\rowcolors{1}{white}{white}
\begin{tabularx}{\textwidth}{>{\centering\arraybackslash}X >{\centering\arraybackslash}X}
\toprule
\rowcolor{gray!25}\textbf{Lemma} & \textbf{Statement} \\
\midrule
\rowcolors{1}{white}{gray!10}
\texttt{pow\_mod3\_q2\_even} & \texttt{q \% 3 = 2 -\textgreater{} q\^{}(2e) \% 3 = 1} --- since q$^{2}$ \equiv 1 (mod 3) \\
\texttt{pow\_mod3\_q2\_odd} & \texttt{q \% 3 = 2 -\textgreater{} q\^{}(2e+1) \% 3 = 2} \\
\texttt{geom\_sum\_mod3\_q2} & \texttt{(\sum\_\{i\textless{}n\} q\^{}i) \% 3 = n \% 2} when \texttt{q \equiv 2 (mod 3)} --- pairs cancel mod 3 \\
\texttt{pow\_mod3\_one} & \texttt{p \% 3 = 1 -\textgreater{} p\^{}k \% 3 = 1} \\
\texttt{sq\_mod3\_of\_not\_dvd} & \texttt{not 3 ∣ m -\textgreater{} m\^{}2 \% 3 = 1} \\
\texttt{sigma\_dvd3\_of\_p2\_kodd} & \texttt{p \% 3 = 2, k \% 2 = 1 -\textgreater{} 3 ∣ sigma(pᵏ)} \\
\bottomrule
\end{tabularx}
\end{table}

\begin{center}
\rule{0.5\textwidth}{0.4pt}
\end{center}

\subsubsection{Main theorems}

\textbf{\texttt{euler\_opn\_form}} \emph{(sorry --- Euler 1747, not yet in Mathlib)}

\begin{quote}
Every OPN has the form \texttt{n = pᵏ . m\^{}2} with \texttt{p} prime, \texttt{p \equiv k \equiv 1 (mod 4)}, \texttt{p ∤ m}.

\end{quote}

\textbf{\texttt{opn\_product\_constraint}}  \emph{fully proved}

\begin{quote}
\end{quote}

\textbf{\texttt{sigma\_prime\_pow\_ratio}}  \emph{fully proved}

\begin{quote}
\end{quote}

\textbf{\texttt{v2\_sigma\_prime\_power}}  \emph{fully proved}

\begin{quote}
For odd prime \texttt{p \equiv 1 (mod 4)} and \texttt{k \equiv 1 (mod 4)}:
\texttt{v₂(sigma(pᵏ)) = 1}

Proof chain: \texttt{sigma(pᵏ) = \sum\_\{i\textless{}=k\} pⁱ} (via \texttt{Finset.sum\_map} after \texttt{Nat.divisors\_prime\_pow}) $\rightarrow$
sum \texttt{\equiv k+1 \equiv 2 (mod 4)} (no LTE needed) $\rightarrow$ \texttt{v₂ = 1}.

\end{quote}

\textbf{\texttt{v2\_sigma\_square\_factor}}  \emph{fully proved}

\begin{quote}
For odd prime \texttt{q} and any \texttt{e}:
\texttt{v₂(sigma(q\^{}(2e))) = 0}

Proof chain: \texttt{sigma(q\^{}(2e)) = \sum\_\{i\textless{}=2e\} qⁱ} $\rightarrow$ sum of \texttt{2e+1} odd terms is odd $\rightarrow$ \texttt{v₂ = 0}.

\end{quote}

\textbf{\texttt{v2\_accumulation\_constraint}}  \emph{fully proved}

\begin{quote}
Given the Euler decomposition \texttt{n = pᵏ . m\^{}2}:
\texttt{v₂(sigma(pᵏ)) = 1}  and  \texttt{forall q \textbar{} m, exists e, v₂(sigma(q\^{}(2e))) = 0}

Machine-verified statement that the 2-adic constraint is \textbf{necessary} for any OPN.
Does not prove nonexistence --- it characterizes what any OPN must look like.

\end{quote}

\textbf{\texttt{opn\_mod4}}  \emph{fully proved}

\begin{quote}
\end{quote}

\textbf{\texttt{touchard\_congruence}}  \emph{fully proved}

\begin{quote}
\end{quote}

\textbf{\texttt{opn\_nonexistence}} \emph{(sorry --- open problem)}

\begin{quote}
\end{quote}

\begin{center}
\rule{0.5\textwidth}{0.4pt}
\end{center}

\subsection{\texttt{Core.lean} --- Primitive type system}

Defines the 12 SynthOmnicon primitives as inductive types with \texttt{deriving DecidableEq, Repr, Ord},
plus explicit \texttt{LE} instances and four cross-primitive axioms:

\begin{table}[htbp]
\centering
\small
\rowcolors{1}{white}{white}
\begin{tabularx}{\textwidth}{>{\centering\arraybackslash}X >{\centering\arraybackslash}X}
\toprule
\rowcolor{gray!25}\textbf{Axiom} & \textbf{Content} \\
\midrule
\rowcolors{1}{white}{gray!10}
A & \texttt{H\_inf -\textgreater{} K\_trap} (topological chirality implies kinetic trapping) \\
B & \texttt{\Omega \textgreater{}= \Omega\_Z -\textgreater{} H \textgreater{}= H2} (integer winding number requires persistent chirality) \\
C & \texttt{D\_holo \textless{}-\textgreater{} T\_holo} (holographic dimensionality iff holographic topology) \\
D & \texttt{\Omega\_NA -\textgreater{} D\_holo} (non-Abelian anyonic protection requires holographic substrate) \\
\bottomrule
\end{tabularx}
\end{table}

\begin{center}
\rule{0.5\textwidth}{0.4pt}
\end{center}

\subsection{Build}

\begin{lstlisting}[language=bash]
lake build SynthOmnicon
\end{lstlisting}

Expected output: \texttt{Build completed successfully} with warnings only (two honest \texttt{sorry} markers,
three unused variable hints from over-specified hypotheses).

\begin{center}
\rule{0.5\textwidth}{0.4pt}
\end{center}

\subsection{Proof-engineering notes}

Several Lean 4.28.0 / Mathlib API subtleties encountered and resolved:

\begin{itemize}
    \item \texttt{Nat.divisors\_prime\_pow} returns \texttt{Finset.map} (with a \texttt{Function.Embedding}), \textbf{not}
\texttt{Finset.image}. Use \texttt{Finset.sum\_map} to unfold, not \texttt{Finset.sum\_image}.
    \item \texttt{omega} cannot cross \texttt{Finset.sum} barriers. Fix: introduce intermediate modular arithmetic
steps that let omega work on plain \texttt{ℕ} expressions.
    \item The induction hypothesis in a \texttt{suffices h : P n by ...} proof carries \texttt{P n}'s
own hypotheses into the IH. Factor out the unconditional general lemma first, then
apply it with the specific hypothesis.
    \item \texttt{norm\_num} primality extension requires \texttt{import Mathlib.Tactic} (not just targeted imports).
    \item \texttt{rw {[}pow\_one{]}} fails inside a \texttt{Finset.sum} after certain rewrite sequences; \texttt{simp} after
\texttt{Finset.sum\_map} handles the residual \texttt{(pⁱ)\^{}1 = pⁱ} correctly.
    \item Numeric \texttt{rw} chains like \texttt{rw {[}Nat.mul\_mod, hq{]}} can fully close goals via \texttt{rfl} when
the resulting expression is a closed numeral (e.g. \texttt{2 * 2 \% 3}). Appending \texttt{norm\_num}
after such a chain produces a ''no goals'' error --- omit it when the \texttt{rw} already closes.
    \item \texttt{absurd h hp3} fails when \texttt{h : 3 = p} but \texttt{hp3 : p != 3} (wrong \texttt{Ne} orientation).
Use \texttt{omega} or \texttt{Ne.symm} to bridge the gap.
    \item \texttt{Dvd.dvd.mul\_left} does not exist in this Mathlib version. Use
\texttt{dvd\_mul\_of\_dvd\_right (dvd\_pow h hn) \_} instead.
\end{itemize}

\begin{center}
\rule{0.5\textwidth}{0.4pt}
\end{center}

\subsection{Design notes and Q\&A}

\subsubsection{Relationship between \texttt{Core.lean} and \texttt{OPN\_2adic.lean}}

\texttt{OPN\_2adic.lean} does not import \texttt{Core.lean}, and at present there is no Lean-level integration
between the two files. The connection is \textbf{conceptual and meta-level}: the SynthOmnicon
primitive framework is used as a thinking tool to organize and motivate the number-theoretic
argument, not as formal input to the proofs.

The specific framing that guided this work: any OPN must resemble a molecule with a unique
''functional group'' (the Euler prime \texttt{pᵏ}, carrying exactly one unit of 2-adic charge) bonded to
an inert ''molecular scaffold'' (the square factor \texttt{m\^{}2}, contributing zero to the 2-adic budget).
The global ''stoichiometric constraint'' \texttt{sigma(n) = 2n} together with the neutrality condition \texttt{n odd}
then severely overdetermines the system. This is the SynthOmnicon constraint-propagation picture
applied to number theory. The primitives that map most naturally onto this picture are:

\begin{table}[htbp]
\centering
\small
\rowcolors{1}{white}{white}
\begin{tabularx}{\textwidth}{>{\centering\arraybackslash}X >{\centering\arraybackslash}X}
\toprule
\rowcolor{gray!25}\textbf{OPN concept} & \textbf{SynthOmnicon analogue} \\
\midrule
\rowcolors{1}{white}{gray!10}
Euler prime \texttt{pᵏ} --- unique 2-adic carrier & \texttt{Phi\_c} --- absorbing under meet, unique criticality carrier \\
Square factors \texttt{q\^{}(2e)} --- 2-adically inert & \texttt{Phi\_sub} with \texttt{K\_trap} isolation \\
\texttt{sigma(n) = 2n} --- global balance & \texttt{S} (Stoichiometry) --- exact ratio constraint \\
\texttt{n} odd --- parity conservation & \texttt{P} (Polarity) --- parity / neutrality condition \\
\bottomrule
\end{tabularx}
\end{table}

The phrase ''mapping down from a more fundamental constraint space'' refers to this conceptual
translation, not to any automated or formal machinery. No custom tactic, external solver, or
category-theoretic functor is involved. The mapping is a heuristic that shaped \emph{which} lemmas
to prove and in what order; correctness is verified entirely by Lean in the usual way.

A formal integration --- where OPN variables are assigned primitive tuples and the constraint
propagation is machine-checked at the SynthOmnicon level --- is a long-term goal but is not yet
planned for Lean.

\begin{center}
\rule{0.5\textwidth}{0.4pt}
\end{center}

\subsubsection{Cross-primitive axioms in \texttt{Core.lean}}

The four \texttt{axiom} declarations are \textbf{foundational postulates} of the SynthOmnicon framework,
not goals to be derived from the inductive type definitions. Within the current Lean
representation --- where each primitive is just a finite inductive type with \texttt{Ord} --- there is
no path to deriving, say, \texttt{D\_holo\_iff\_T\_holo} from first principles, because the types carry
no algebraic structure that encodes the physical relationship between dimensionality and topology.
Deriving such equivalences would require an external interpretation: a functor from the primitive
lattice into a category where boundary/bulk duality has an intrinsic meaning (e.g., a
topos-theoretic or operadic model). That is out of scope for the current formalization.

\textbf{On \texttt{LE} instances vs \texttt{LinearOrder}:} the custom \texttt{LE} via \texttt{compare} is intentionally minimal.
The types that need it (\texttt{Protection}, \texttt{Chirality}) only use \texttt{\textgreater{}=} in the cross-primitive axioms,
not lattice operations. Deriving a full \texttt{LinearOrder} is straightforward --- all types derive \texttt{Ord}
and the machinery exists in Mathlib --- but premature: it would be needed only once \texttt{⊓}/\texttt{⊔}
operations are actually used in proofs. The exception is \texttt{Criticality}, where \texttt{Phi\_c} is
\emph{absorbing} under meet (\texttt{meet(Phi\_c, x) = Phi\_c} for all \texttt{x}), which contradicts \texttt{min} semantics.
A correct \texttt{MeetSemilattice} instance for \texttt{Criticality} will require a custom definition that
overrides the \texttt{Ord}-derived ordering for meet --- this is noted in a comment in \texttt{Core.lean} and
is the first planned extension to the primitive lattice.

\begin{center}
\rule{0.5\textwidth}{0.4pt}
\end{center}

\subsubsection{OPN formalization details}

\textbf{Coprimality hypothesis in \texttt{touchard\_congruence}.} The \texttt{hcop : Nat.Coprime (p\^{}k) (m\^{}2)} in
\texttt{touchard\_congruence} (and in \texttt{opn\_product\_constraint}) is an explicit parameter rather than
something derived, because \texttt{euler\_opn\_form} is still \texttt{sorry}. Once \texttt{euler\_opn\_form} is proved,
\texttt{hcop} becomes derivable: \texttt{p} prime, \texttt{p ∤ m} $\rightarrow$ \texttt{Nat.Coprime p m} $\rightarrow$
\texttt{Nat.Coprime (p\^{}k) (m\^{}2)} by \texttt{Nat.Coprime.pow}. At that point, \texttt{hcop} can be dropped from the
signature and threaded through internally. The current signature is a faithful statement of what
the proof actually uses, without hiding a dependency on the sorry.

\textbf{\texttt{factorization.support} vs \texttt{forall q, q.Prime -\textgreater{} q ∣ m}.} The \texttt{factorization.support} indexing is
intentional: it provides the exponent \texttt{(Nat.factorization m) q} directly, which is exactly the
\texttt{e} needed for the existential \texttt{exists e, v₂(sigma(q\^{}(2e))) = 0}. An alternative phrasing with
\texttt{q.Prime -\textgreater{} q ∣ m} would require a separate argument that the correct exponent exists and is
given by the factorization --- more work for no gain in readability.

\textbf{\texttt{euler\_opn\_form} in Mathlib.} As of Lean 4.28.0 / Mathlib v4.28.0, the Euler form for OPNs
(\texttt{n = pᵏm\^{}2}, \texttt{p \equiv k \equiv 1 mod 4}) is \textbf{not in Mathlib}. The number theory coverage in Mathlib
is extensive for multiplicative functions, divisor sums, and primality, but this specific
structural theorem --- which requires a careful argument about the 2-adic structure of \texttt{sigma(n) = 2n}
applied globally across the full factorization --- has not been formalized. Proving it from the
tools already in this file is feasible: the key lemmas (\texttt{sigma\_mul\_of\_coprime},
\texttt{sigma\_prime\_pow\_ratio}, \texttt{v2\_eq\_one\_of\_mod4\_eq2}) are all in place. It is a planned extension,
not a fundamental obstacle.

\begin{center}
\rule{0.5\textwidth}{0.4pt}
\end{center}

\subsubsection{Other Mathlib API surprises}

Beyond the items already in the proof-engineering notes:

\begin{itemize}
    \item \textbf{\texttt{zify} is required to use \texttt{geom\_sum\_mul}.} The Mathlib \texttt{geom\_sum\_mul} lemma operates
over a \texttt{CommRing}, so it lives in ℤ (or any ring), not ℕ. The \texttt{zify} tactic --- with guards
\texttt{{[}hp.one\_le, Nat.one\_le\_pow ...{]}} to handle the coercion side conditions --- is the cleanest
bridge. Attempting to prove the geometric identity directly in ℕ avoids the sign issues but
requires manual subtraction accounting.
    \item \textbf{\texttt{v2\_eq\_one\_of\_mod4\_eq2} design choice.} The \texttt{by\_contra} / \texttt{omega} approach was preferred
over a more direct use of \texttt{Prime.pow\_dvd\_iff\_le\_factorization} because the contrapositive
direction (\texttt{factorization \textgreater{}= 2 -\textgreater{} 2\^{}2 ∣ n}) is what the lemma naturally produces, and making
both the lower and upper bound explicit as separate steps keeps the proof readable. A more
compressed proof using \texttt{Nat.le\_antisymm} with two applications of
\texttt{pow\_dvd\_iff\_le\_factorization} is possible and would be slightly shorter; the current form
was chosen for clarity.
    \item \textbf{Unused variable warnings} in \texttt{v2\_accumulation\_constraint} (\texttt{h\_perf}) and
\texttt{v2\_sigma\_prime\_power} (\texttt{hp\_odd}) are intentional: those hypotheses are included because
they are part of the natural statement of the theorem (an odd perfect number has an odd prime
in its Euler form) even though the specific proof path doesn't reach them. They document
the full precondition rather than minimizing the signature. They can be removed without
affecting correctness.
\end{itemize}

\begin{center}
\rule{0.5\textwidth}{0.4pt}
\end{center}

\subsubsection{The ''mapping down'' claim --- honest account}

There is no automated reasoning, custom tactic, or external solver involved. The phrase
''mapping down from a more fundamental constraint space'' refers to using the SynthOmnicon
conceptual framework --- specifically its vocabulary of unique constraint carriers, neutral
scaffolds, and overdetermined global balance equations --- as a \emph{problem-structuring heuristic}.
This heuristic suggested:

\begin{enumerate}
    \item Look for the unique carrier of the conserved quantity (v₂ budget = 1, carried by \texttt{pᵏ}).
    \item Verify the scaffold is transparent (square factors contribute zero).
    \item Write down the overdetermined global constraint and check consistency.
    \item Apply the same template at the next prime (3-adic version for Touchard).
\end{enumerate}

All proofs are standard Lean 4 / Mathlib proofs verified by the kernel. The framework provided
direction, not machinery.

\begin{center}
\rule{0.5\textwidth}{0.4pt}
\end{center}

\subsubsection{Future directions}

\textbf{Connecting the two tracks.} The most concrete near-term connection: as \texttt{Core.lean} gains
\texttt{Semilattice} / \texttt{Lattice} instances, the OPN constraint structure could be encoded as a
\emph{synthon} --- a tuple in the 12-dimensional primitive space --- and the constraint propagation
verified at that level. The \texttt{Phi\_c} absorbing-meet property is the structural analogue of
the Euler prime's uniqueness; formalizing this analogy in Lean would be the first real
integration between the files. This is planned but not yet scoped.

\textbf{Lattice instances for \texttt{Core.lean}.} The immediate next step for \texttt{Core.lean} is implementing
the custom \texttt{MeetSemilattice} for \texttt{Criticality} (where \texttt{Phi\_c} is absorbing). After that, full
\texttt{Lattice} instances for the five ordered primitives (\texttt{F}, \texttt{K}, \texttt{G}, \texttt{\Omega}, \texttt{H}) are
straightforward using \texttt{minFac}/\texttt{maxFac} over the \texttt{Ord}-derived linear order. The categorical
primitives (\texttt{D}, \texttt{T}, \texttt{R}, \texttt{P}, \texttt{\Gamma}) require a different treatment: their meet semantics are
identity-or-bottom, not min, so they need a custom \texttt{BoundedLattice} with \texttt{\perp} representing
''incompatible''.

\textbf{OPN track.} The natural next steps in order of difficulty:

\begin{enumerate}
    \item Prove \texttt{euler\_opn\_form} (all required tools are present).
    \item Extend to \texttt{touchard\_congruence\_v2} without the Euler decomposition as a hypothesis
(i.e., derive the decomposition inside the proof from \texttt{euler\_opn\_form}).
    \item Begin the prime bound argument: from \texttt{opn\_product\_constraint} and \texttt{sigma\_prime\_pow\_ratio},
derive a lower bound on the number of distinct prime factors of any OPN.
\end{enumerate}


\end{document}
